\subsection{张量与张量记号}

设 A 是一个交换环，M 是一个 A-模， \(\mathbf{M}^*\) 为 M 的对偶（II，§2，第3号），I 和 J 为两个不相交的有限集；A-模 \(\bigotimes_{i\in \mathrm{I}\cup \mathrm{J}}\mathbf{E}_i\) （其中 \(\mathbf{E}_i = \mathbf{M}\) （若 \(i\in \mathrm{I}\) ）， \(\mathbf{E}_i = \mathbf{M}^*\) （若 \(i\in \mathrm{J}\) ））记作 \(\mathsf{T}_{\mathsf{J}}^{\mathsf{I}}(\mathbf{M})\) ； \(\mathsf{T}_{\mathsf{J}}^{\mathsf{I}}(\mathbf{M})\) 的元素称为 M 上的 (I,J) 型张量。当 \(\mathbf{J} = \varnothing\) 时它们称为反变张量，当 \(\mathbf{I} = \varnothing\) 时称为协变张量，其余情形称为混合张量。

设 \(I'\) , \(I''\) 为 I 的两个子集，且 \(J'\) , \(J''\) 为 J 的两个子集，使得 \(I' \cup I'' = I\) , \(\mathbf{I}' \cap \mathbf{I}'' = \varnothing, \mathbf{J}' \cup \mathbf{J}'' = \mathbf{J}, \mathbf{J}' \cap \mathbf{J}'' = \varnothing\) ；则存在一个典范的结合同构（II，§3，第9号）

\[
\mathsf {T} _ {\mathbf {J}} ^ {\mathbf {I}} (\mathbf {M}) \rightarrow \mathsf {T} _ {\mathbf {J} ^ {\prime}} ^ {\mathbf {I} ^ {\prime}} (\mathbf {M}) \otimes_ {\mathbf {A}} \mathsf {T} _ {\mathbf {J} ^ {\prime \prime}} ^ {\mathbf {I} ^ {\prime \prime}} (\mathbf {M}).\tag{11}
\]

考虑张量代数 \(\mathsf{T}(\mathbf{M}\oplus \mathbf{M}^{*})\) ，由第 5 号可知 \(\mathsf{T}^n (\mathbf{M}\oplus \mathbf{M}^*)\) 典范地等同于诸 \(\mathsf{T}_{\mathbf{J}}^{\mathbf{I}}(\mathbf{M})\) 的直和，其中 I 遍历区间 \([1,n]\) （在 \(\mathbf{N}\) 中）的子集组成的集合，而 J 是 I 在 \([1,n]\) 中的补集。

当 \(I = [1, p]\) 且 \(J = [p + 1, p + q]\) ，其中整数 \(p \geqslant 0\) , \(q \geqslant 0\) （我们将 I（相应地 J）替换为 \(\varnothing\) ，如果 \(p = 0\) （相应地 \(q = 0\) ）），A-模 \(T_J^I(M)\) 也记作 \(T_q^p(M)\) ；A-模 \(T_0^n(M)\) 和 \(T_n^0(M)\) 根据定义分别是 A-模 \(T^n(M)\) 和 \(T^n(M^*)\) 。当 I 和 J 是基数为 \(p = \text{Card}(I)\) 和 \(q = \text{Card}(J)\) 的任意有限集合时，我们给它们每一个赋予一个全序；那么存在一个从 I（相应地 J）到 \([1, p]\) （相应地 \([p + 1, p + q]\) ）上的递增双射，因此这些双射定义了一个同构

\[
\mathbf {T} _ {\mathrm{J}} ^ {\mathrm{I}} (\mathbf {M}) \rightarrow \mathbf {T} _ {q} ^ {p} (\mathbf {M}).
\]

当 M 是有限生成的投射 A-模时，由 II，§ 2，第 7 号，命题 13 的推论 4 以及 II，§ 4，第 4 号，命题 4 的推论 1 可知，存在一个典范同构

\[
(\mathsf {T} _ {\mathbf {J}} ^ {\mathbf {I}} (\mathbf {M})) ^ {*} \to \mathsf {T} _ {\mathbf {I}} ^ {\mathbf {J}} (\mathbf {M}).
\]

现在假设 \(\mathbf{M}\) 是一个有限生成的自由 A-模，且设 \((e_{\lambda})_{\lambda \in L}\) 为 \(\mathbf{M}\) 的一组基（因此 L 是有限集）。 \(\mathbf{M}^*\) 的与 \((e_{\lambda})\) 对偶的基（II，§2，第 6 号）记为 \((e^{\lambda})_{\lambda \in L}\) 。基 \((e_{\lambda})\) 和 \((e^{\lambda})\) （分别属于 \(\mathbf{M}\) 和 \(\mathbf{M}^*\) ）定义了（第 5 号） \(\mathsf{T}_{\mathbf{J}}^{\mathbf{I}}(\mathbf{M})\) 的一个基，我们明确给出如下：给定两个映射 \(f\colon \mathbf{I} \to \mathbf{L}\) 和 \(g\colon \mathbf{J} \to \mathbf{L}\) ，设 \(e_f^g\) 为元素 \(\bigotimes_{i \in I \cup J} x_i\) （属于 \(\mathsf{T}_{\mathbf{J}}^{\mathbf{I}}(\mathbf{M})\) ），定义为

\[
x _ {i} = e _ {f (i)} \quad \text { if } i \in \mathbf {I}, \qquad x _ {i} = e ^ {g (i)} \quad \text { if } i \in \mathbf {J}.
\]

当 \((f, g)\) 遍历映射 \(f: \mathbf{I} \to \mathbf{L}\) 和 \(g: \mathbf{J} \to \mathbf{L}\) 的有序对组成的集合时，诸 \(e_f^g\) 构成 A-模 \(\mathsf{T}_{\mathbf{J}}^{\mathbf{I}}(\mathbf{M})\) 的一组基，称为与给定基 \((e_{\lambda})\) （ \(\mathbf{M}\) 的）相关联的基。对于 \(z \in \mathsf{T}_{\mathbf{J}}^{\mathbf{I}}(\mathbf{M})\) ，我们因此可以写出

\[
z = \sum_ {(f, g)} \alpha_ {g} ^ {f} (z). e _ {f} ^ {g}\tag{12}
\]

其中 \(\alpha_g^f\) 是关于基 \((e_f^g)\) 的坐标线性型；按照语言的约定俗成， \(\alpha_g^f(z)\) 称为张量 \(z\) 关于模 M 的基 \((e_\lambda)\) 的坐标。诸 \(\alpha_g^f\) 构成 \((e_g^f)\) 的对偶基，换言之，它们与 \(\mathsf{T}_{\mathbf{I}}^{\mathbf{j}}(\mathbf{M})\) 的基的元素等同，该基与 \((e_\lambda)\) 相关联。当 I 和 J 是 N 中区间 \([1,n]\) 的互补子集时， \(\alpha_g^f\) （或 \(\alpha_g^f(z)\) ）记法如下：每个元素 \(f(i)\) （对于 \(i\in I\) ）写作第 \(i\) 位上的上标，并在第 \(i\) 位处放一个点（对 \(i \in \mathbf{J}\) ）；类似地， \(g(i)\) （对于 \(i \in \mathbf{J}\) ）写作第 \(i\) 位上的下标，并在第 \(i\) 位处放一个点（对 \(i \in \mathbf{I}\) ）。例如，对于 \(\mathbf{I} = \{1, 4\}\) , \(\mathbf{J} = \{2, 3\}\) ，我们记 \(\alpha_{\cdot \nu \rho}^{\lambda \cdot \cdot \mu}\) ，如果 \(f(1) = \lambda, f(4) = \mu, g(2) = \nu, g(3) = \rho\) 。

设 \((\bar{e}_{\lambda})_{\lambda \in L}\) 为 M 的另一组基，且 \(P\) 为从 \((e_{\lambda})\) 到 \((\bar{e}_{\lambda})\) 的过渡矩阵（II，§ 10，第 8 号）。则从 \((e^{\lambda})\) 到 \((\bar{e}^{\lambda})\) （ \((\bar{e}_{\lambda}))\) 的对偶基）的过渡矩阵是 \(^{t}P^{-1}\) ，即 \(P\) 的逆步（II，§ 10，第 8 号，命题 5）。由此可得（II，§ 10，第 10 号），从基 \((e_{f}^{g})\) （属于 \(\mathrm{T}_{\mathrm{J}}^{\mathrm{I}}(\mathrm{M})\) ）到基 \((\bar{e}_{f}^{g})\) (其中 \(f\) （相应地 \(g\) ) 遍历从 I 到 L（相应地，从 J 到 L）的映射集合）的过渡矩阵是矩阵

\[
\bigotimes_ {i \in I \cup J} Q _ {i}, \quad \text { where } Q _ {i} = P \text { if } i \in I, Q _ {i} = ^ {t} P ^ {- 1} \text { if } i \in J.\tag{13}
\]

因此，该矩阵的转置被等同于

\[
\bigotimes_ {i \in \mathbf {I} \cup J} R _ {i}, \quad \text { where } R _ {i} = ^ {t} P ^ {- 1} \text { if } i \in \mathbf {I}, R _ {i} = P \text { if } i \in \mathbf {J}.\tag{14}
\]

现在假设 \(\mathbf{M}\) 是任意模。设 \(i\in \mathbf{I}, j\in \mathbf{J}\) 并记 \(\mathbf{I}' = \mathbf{I} - \{i\}, J' = J - \{j\}\) ；我们将定义一个典范的 A-线性映射

\[
c _ {j} ^ {i} \colon \mathsf {T} _ {\mathbf {J}} ^ {\mathbf {I}} (\mathbf {M}) \to \mathsf {T} _ {\mathbf {J} ^ {\prime}} ^ {\mathbf {I} ^ {\prime}} (\mathbf {M}),
\]

称为指标 \(i\) 与指标 \(j\) 的缩并。为此，注意到 \(\mathbf{M}^{\mathrm{I}} \times (\mathbf{M}^{*})^{\mathrm{J}}\) 上的映射，它将每个族 \((x_{i})_{i \in \mathbf{I} \cup \mathbf{J}}\) （其中 \(x_{i} \in \mathbf{M}\) （若 \(i \in \mathbf{I}\) ）， \(x_{i} \in \mathbf{M}^{*}\) （若 \(i \in \mathbf{J}\) ））对应到元素

\[
\left\langle x _ {i}, x _ {j} \right\rangle \otimes_ {k \in (\mathbb {I} \cup \mathbb {J}) - \{i, j \}} x _ {k}\tag{15}
\]

（属于 \(T_{J'}^{I'}(M)\) ）是 A-多重线性的； \(c_{j}^{i}\) 是对应的 A-线性映射。

现在假设 M 是自由且有限生成的，并令 \((e_{\lambda})_{\lambda \in L}\) 为 M 的一组基。给定两个映射 \(f: I \to L\) , \(g: J \to L\) ，设 \(f_i\) 表示 f 在 \(I' = I - \{i\}\) 上的限制， \(g_j\) 表示 g 在 \(J' = J - \{j\}\) 上的限制；于是根据 (12)

\[
c _ {j} ^ {i} (e _ {f} ^ {g}) = \left\{ \begin{array}{l l} 0 & \text { if } f (i) \neq g (j) \\ e _ {f _ {i}} ^ {g _ {j}} & \text { if } f (i) = g (j) \end{array} \right.\tag{16}
\]

\(c_j^i(z)\) 的坐标作为 \(z\) 的坐标的函数的表达式由此得出；对于每个映射 \(f'\) （相应地 \(g'\) ）——从 \(I'\) 到 \(L\) （相应地，从 \(J'\) 到 \(L\) ）——以及每个 \(\lambda \in L\) ，设 \((f', \lambda)\) （相应地 \((g', \lambda)\) ）表示从 \(I\) 到 \(L\) （相应地，从 \(J\) 到 \(L\) ）的、在 \(I'\) （相应地 \(J'\) ）上的限制为 \(f'\) （相应地 \(g'\) ）、并且以 \(\lambda\) 为其在元素 \(i\) （相应地 \(j\) ）处的值的映射。那么，若关于基 \((e_{f'}^{g'})\) （ \(T_{J'}^{I'}(M)\) 的基）的坐标线性型记为 \(\beta_{g'}^{f'}\) ,

\[
\beta_ {g ^ {\prime}} ^ {f ^ {\prime}} (c _ {j} ^ {i} (z)) = \sum_ {\lambda \in L} \alpha_ {(g ^ {\prime}, \lambda)} ^ {(f ^ {\prime}, \lambda)} (z).\tag{17}
\]

张量的例子。(1) 设 M 是有限生成的投射 A-模。我们知道（II，§ 4，第 2 号，命题 2 的推论），存在一个典范的 A-模同构

\[
\theta_ {\mathbf {M}} \colon \mathbf {M} ^ {*} \otimes_ {\mathbf {A}} \mathbf {M} \rightarrow \operatorname{End} _ {\mathbf {A}} (\mathbf {M})
\]

，使得 \(\theta_{\mathbf{M}}(x^{*}\otimes x)\) （对于 \(x\in \mathbf{M},x^{*}\in \mathbf{M}^{*}\) ）是自同态

\[
y \mapsto \langle y, x ^ {*} \rangle x.
\]

因此，借助 \(\theta_{\mathbf{M}}\) , \(\mathsf{T}_{(1)}^{(2)}(\mathbf{M})\) （同构于 \(\mathsf{T}_1^1 (\mathbf{M}))\) 可与 A-模 \(\operatorname{End}_{\mathbf{A}}(\mathbf{M})\) 等同。假设 \(\mathbf{M}\) 是自由模，且设 \((e_{\lambda})_{\lambda \in \mathbf{L}}\) 为 \(\mathbf{M}\) 的一组基；那么张量 \(z\in \mathbf{M}^*\otimes \mathbf{M}\) 关于该模的基 \((e^{\mu}\otimes e_{\lambda})\) 的坐标记为 \(\zeta_{\mu}^{\cdot \lambda}\) 。由于 \(\theta_{\mathrm{M}}(e^{\mu}\otimes e_{\lambda})\) 是自同态 \(y\mapsto \langle y,e^{\mu}\rangle e_{\lambda}\) ，自同态 \(u = \theta_{\mathrm{M}}(z) = \theta_{\mathrm{M}}\left(\sum_{\lambda ,\mu}\zeta_{\mu}^{\cdot \lambda}e^{\mu}\otimes e_{\lambda}\right)\) 把 \(y\) 映到 \(\sum_{\lambda ,\mu}\zeta_{\mu}^{\cdot \lambda}\langle y,e^{\mu}\rangle e_{\lambda}\) ；取 \(y = e_{\lambda}\) ，我们得到关系式

\[
u (e _ {\lambda}) = \sum_ {\mu \in L} \zeta_ {\lambda} ^ {\cdot \mu} e _ {\mu}\tag{18}
\]

换言之，线性映射 \(u = \theta_{\mathrm{M}}(z)\) 的矩阵是这样的矩阵：其位于指标 \(\mu\) 的行与指标 \(\lambda\) 的列交叉处的元素为 \(\zeta_{\lambda}^{\mu}\) 。

\(u\) 的迹的定义（II，§4，第3号）立即表明

\[
\operatorname{Tr} (\theta_ {\mathbb {M}} (z)) = c _ {1} ^ {2} (z).
\]

因此，元素 \(z_0 = \sum_{\lambda \in L} e^\lambda \otimes e_\lambda\) （其坐标 \(\zeta_\mu^{\cdot \lambda}\) 对于 \(\lambda \neq \mu\) 为零，对于 \(\lambda = \mu\) 等于1），它满足 \(\theta_M(z_0) = 1_M\) ，是元素 \(1 \in A = T_0^0(M)\) 在映射（缩并的转置）下的像。

\[
c _ {1} ^ {2}: \mathbf {T} _ {\{1 \}} ^ {\{2 \}} (\mathbf {M}) \rightarrow \mathbf {A}.
\]

\begin{enumerate}
\def\labelenumi{(\arabic{enumi})}
\setcounter{enumi}{1}
\tightlist
\item
  始终假设 M 是有限生成的投射 A-模；存在一个典范的 A-模同构
\end{enumerate}

\[
\mu : \mathbf {M} ^ {*} \otimes_ {\mathbf {A}} \mathbf {M} ^ {*} \rightarrow (\mathbf {M} \otimes_ {\mathbf {A}} \mathbf {M}) ^ {*}
\]

（II，§4，第 4 号，命题 2 的推论 1）以及一个典范同构

\[
\theta : (\mathbf {M} \otimes_ {\mathbf {A}} \mathbf {M}) ^ {*} \otimes_ {\mathbf {A}} \mathbf {M} \rightarrow \operatorname{Hom} _ {\mathbf {A}} (\mathbf {M} \otimes_ {\mathbf {A}} \mathbf {M}, \mathbf {M})
\]

（II，§4，第 2 号，命题 2 的推论）；此外 \(\operatorname{Hom}_{\mathbf{A}}(\mathbf{M} \otimes_{\mathbf{A}} \mathbf{M}, \mathbf{M})\) 典范同构于 A-模 \(\mathcal{L}_2(\mathbf{M}, \mathbf{M}; \mathbf{M})\) ，即由从 \(\mathbf{M} \times \mathbf{M}\) 到 \(\mathbf{M}\) 的 A-双线性映射组成的模（II，§3，第 9 号）。复合这些同构，得到典范同构

\[
\chi_ {\mathbf {M}} \colon T _ {\{1, 2 \}} ^ {(3)} (\mathbf {M}) = \mathbf {M} ^ {*} \otimes \mathbf {M} ^ {*} \otimes \mathbf {M} \rightarrow \mathcal {L} _ {2} (\mathbf {M}, \mathbf {M}; \mathbf {M})
\]

使得，对于 \(x^{*}, y^{*}\) （属于 \(\mathbf{M}^{*}\) ）， \(z \in \mathbf{M}\) ， \(\chi_{\mathbf{M}}(x^{*} \otimes y^{*} \otimes z)\) 是双线性映射

\[
(u, v) \mapsto \langle u, x ^ {*} \rangle \langle v, y ^ {*} \rangle z.
\]

因此，借助 \(\chi_{\mathbf{M}}\) , \(\mathsf{T}_{\{1,2\}}^{(3)}(\mathbf{M})\) （同构于 \(\mathsf{T}_{2}^{1}(\mathbf{M})\) ）可与 A-模 \(\mathcal{L}_2(\mathbf{M},\mathbf{M};\mathbf{M})\) 等同。假设 \(\mathbf{M}\) 是一个自由 A-模，且设 \((e_{\lambda})_{\lambda \in L}\) 为 \(\mathbf{M}\) 的一组基；那么张量 \(z\in \mathbf{M}^*\otimes \mathbf{M}^*\otimes \mathbf{M}\) 关于该模的基 \((e^{\lambda}\otimes e^{\mu}\otimes e_{\nu})\) 的坐标记为 \(\zeta_{\lambda \mu}^{\cdot \cdot \nu}\) 。双线性映射 \(\chi_{\mathbf{M}}(z)\) 把有序对 \((e_{\lambda},e_{\mu})\) 映到

\[
\sum_ {\nu \in L} \zeta_ {\lambda \mu}^{\cdot \cdot \nu} e _ {\nu}
\]

因此诸 \(\zeta_{\lambda \mu}^{\cdots \nu}\) 正是定义在 \(\mathbf{M}\) 上、由双线性映射 \(\chi_{\mathbf{M}}(z)\) 给出的（一般非结合的）代数关于基 \((e_{\lambda})\) 的结构常数（ \(\S 1\) ，第 7 号）。

注记2. 设 \((e_{\lambda})_{\lambda \in L}, (\tilde{e}_{\lambda})_{\lambda \in L}\) 为 M 的两个基，且 \(P\) 为从 \((e_{\lambda})\) 到 \((\tilde{e}_{\lambda})\) 的过渡矩阵。根据例1中所见， \(P\) 中位于指标 \(\lambda\) 的行与指标 \(\mu\) 的列交叉处的元素记为 \(\alpha_{\mu}^{\lambda}\) ，逆步 \(^{t}P^{-1}\) 中位于指标 \(\lambda\) 的行与指标 \(\mu\) 的列交叉处的元素记为 \(\beta_{\lambda}^{\mu}\) 。于是（采用上面引入的记号）

\begin{enumerate}
\def\labelenumi{(\arabic{enumi})}
\setcounter{enumi}{18}
\tightlist
\item
\end{enumerate}

\[
\left\{ \begin{array}{l} \tilde {e} _ {\mu} = \sum_ {\lambda} \alpha_ {\mu} ^ {\lambda} e _ {\lambda} \\ \tilde {e} ^ {\mu} = \sum_ {\lambda} \beta_ {\lambda} ^ {\mu} e ^ {\lambda} \end{array} \right.\tag{20}
\]

\[
\bar {e} _ {f ^ {\prime}} ^ {g ^ {\prime}} = \sum_ {(f, g)} \left(\prod_ {i \in I} \alpha_ {f ^ {\prime} (i)} ^ {f (i)}\right) \left(\prod_ {j \in J} \beta_ {g (j)} ^ {g ^ {\prime} (j)}\right) e _ {f} ^ {g}
\]

对于所有映射 \(f': \mathbf{I} \to \mathbf{L}\) 和 \(g': \mathbf{J} \to \mathbf{L}\) 成立。于是，坐标 \(\zeta_g^f\) ——张量 \(z \in \mathbb{T}_J^{\mathbf{I}}(\mathbf{M})\) 关于基 \((e_{\lambda})\) 的坐标——可以利用公式由坐标 \(\bar{\zeta}_{g'}^{f'}\) —— \(z\) 关于基 \((\bar{e}_{\lambda})\) 的坐标——表示出来

\[
\zeta_ {g} ^ {f} = \sum_ {(f ^ {\prime}, g ^ {\prime})} \left(\prod_ {i \in I} \alpha_ {f ^ {\prime} (i)} ^ {f (i)}\right) \left(\prod_ {j \in J} \beta_ {g (j)} ^ {g ^ {\prime} (j)}\right) \bar {\zeta} _ {g ^ {\prime}} ^ {f ^ {\prime}}.\tag{21}
\]

矩阵 \(P^{-1}\) （从基 \((\tilde{e}_{\lambda})\) 到基 \((e_{\lambda})\) 的过渡矩阵）是 \(^{t}P^{-1}\) 的转置，使得 \(\beta_{\lambda}^{\mu}\) 是位于指标 \(\lambda\) 的列与指标 \(\mu\) 的行交叉处、出现在 \(P^{-1}\) 中的元素。对 \(e_{f}^{g}\) （用 \(\tilde{e}_{f'}^{g'}\) 表示）以及对 \(\overline{\zeta}_{g'}^{f'}\) （用 \(\zeta_{g}^{f}\) 表示）的计算，因此是通过在上述计算中把 \(\alpha_{\mu}^{\lambda}\) 换为 \(\beta_{\lambda}^{\mu}\) 、把 \(\beta_{\lambda}^{\mu}\) 换为 \(\alpha_{\mu}^{\lambda}\) ，并交换 \(f\) 与 \(f'\) 以及 \(g\) 与 \(g'\) 的角色而得到的。则

\[
e _ {f} ^ {g} = \sum_ {(f ^ {\prime}, g ^ {\prime})} \left(\prod_ {j \in J} \alpha_ {g ^ {\prime} (j)} ^ {g (j)}\right) \left(\prod_ {i \in I} \beta_ {f (i)} ^ {f ^ {\prime} (i)}\right) \bar {e} _ {f ^ {\prime}} ^ {g ^ {\prime}}\tag{22}
\]

\[
\bar {\zeta} _ {g ^ {\prime}} ^ {f ^ {\prime}} = \sum_ {(f, g)} \left(\prod_ {j \in J} \alpha_ {g ^ {\prime} (j)} ^ {g (j)}\right) \left(\prod_ {i \in I} \beta_ {f (i)} ^ {f ^ {\prime} (i)}\right) \zeta_ {g} ^ {f}.\tag{23}
\]

上述公式的特点是求和针对那些分别作为下指标和上指标各出现一次的指标进行。某些作者鉴于这一情况，允许自己省略求和符号。

